% =============================================================================
\chapter{Assembly, Boundary Conditions, and Constraints}
\label{ch:assembly}
% =============================================================================

This chapter connects local element equations to global algebraic systems. It
focuses on local-to-global mappings, constrained degrees of freedom, and the
different mathematical interpretations of essential and natural boundary
conditions.

\section{Element connectivity and global assembly maps}
\label{sec:assembly:connectivity}

% Global numbering, sparsity patterns, and assembly graphs

\section{The assembly of linear and multilinear forms}
\label{sec:assembly:forms}

% Accumulation of vectors, matrices, and higher-order tensors from element contributions

\section{Essential and natural boundary conditions}
\label{sec:assembly:boundary}

The FE layer distinguishes pointwise essential data from trace-essential data.
For standard $H^1$ spaces, strong Dirichlet conditions remain pointwise
constraints on nodal degrees of freedom. For trace-conforming vector spaces,
such as $H(\mathrm{div})$, the essential datum is not a point value but a
scalar trace. The canonical example is the normal trace $u \cdot n$.

This distinction matters because the FE surface must remain physics-agnostic.
The library therefore exposes generic trace-condition objects rather than
Darcy-specific or contact-specific boundary types. Equality-based trace laws
are represented by strong normal-trace constraints, weak trace loads, or weak
Robin relations. Boundary and interface traces use the same vocabulary, so the
same scalar trace operator can be applied on $\mathrm{d}s(\Gamma)$ or on
interior interface measures.

\section{Strong imposition of constraints}
\label{sec:assembly:strong_constraints}

Strong constraints are algebraic statements about DOFs. They are applied
through elimination or lifting and therefore do not alter the variational form
directly. Periodicity, tied-node relations, and strong $H(\mathrm{div})$
normal-trace data all live in this category. The modern FE surface therefore
includes both pointwise Dirichlet constraints and orientation-aware trace
constraints for vector-basis spaces.

\section{Weak imposition of constraints}
\label{sec:assembly:weak_constraints}

Weak constraints remain part of the residual and Jacobian. The FE library now
provides a unified trace-oriented weak-constraint surface:

\begin{itemize}
\item load and Robin trace terms on boundaries and interfaces,
\item trace-oriented Nitsche terms for equality-style weak imposition,
\item one-sided scalar-trace inequalities for semismooth nonlinear solves.
\end{itemize}

This taxonomy avoids encoding PDE-specific semantics in FE while still letting
applications express normal-flux, exchange, penalty, and unilateral laws in a
mathematically consistent way.

\section{Multi-point constraints and nonmatching interfaces}
\label{sec:assembly:mpc}

Multi-point constraints include periodic relations, tied interfaces, and other
algebraic trace relations that do not belong in the weak form. For
orientation-sensitive vector-basis spaces, matching coordinates alone are not
sufficient; the trace relation must also account for facet ordering and outward
normal orientation. The FE helper surface therefore includes dedicated
$H(\mathrm{div})$ periodic and MPC trace utilities.

For nonconforming or hybridized couplings, the FE layer now also supports
mortar-style facet fields. These fields are assembled on interface faces,
occupy their own algebraic block, and couple to neighboring volume fields
through interface kernels rather than through pointwise tied DOF logic.

\section{Structure of the global algebraic system}
\label{sec:assembly:global_system}

Modern FE systems are rarely a single nodal block. Trace fields, mortar
multipliers, and mixed-dimensional interface fields all introduce additional
rows and columns that are topologically lower-dimensional but algebraically
first-class. The global system can therefore contain:

\begin{itemize}
\item volume-volume blocks,
\item volume-interface coupling blocks,
\item interface-interface or mortar-mortar facet blocks,
\item weak trace-law contributions and one-sided nonlinear trace terms.
\end{itemize}

The practical consequence is that assembly metadata, sparsity construction, and
solver strategy must understand these blocks as part of the core FE
infrastructure rather than as afterthoughts layered on a purely cell-based
system.
